\documentclass[a4paper,12pt]{article}
\usepackage[utf8]{inputenc}
\usepackage[ngerman]{babel}
\usepackage{geometry}
\usepackage{graphicx}
\usepackage{amsmath, amssymb}
\usepackage{booktabs}
\usepackage{array}
\usepackage{longtable}
\usepackage{hyperref}
\usepackage{xcolor}
\usepackage{fancyhdr}
\usepackage{titlesec}
\usepackage{abstract}
\usepackage{amsthm}
\usepackage{cite}
\usepackage{microtype}
\usepackage{setspace}
\usepackage{caption}
\usepackage{subcaption}
\usepackage{float}
\usepackage{listings}
\usepackage{enumitem}
\usepackage{tikz}
\usetikzlibrary{positioning}
\usepackage{pgfplots}
\pgfplotsset{compat=1.18}

\definecolor{accentblue}{RGB}{0,90,180}

\lstdefinelanguage{json}{
  basicstyle=\ttfamily\small,
  showstringspaces=false,
  breaklines=true,
  morestring=[b]",
  morecomment=[l]{//},
  morecomment=[s]{/*}{*/},
  literate=
    *{0}{{\textcolor{blue!70!black}{0}}}{1}
     {1}{{\textcolor{blue!70!black}{1}}}{1}
     {2}{{\textcolor{blue!70!black}{2}}}{1}
     {3}{{\textcolor{blue!70!black}{3}}}{1}
     {4}{{\textcolor{blue!70!black}{4}}}{1}
     {5}{{\textcolor{blue!70!black}{5}}}{1}
     {6}{{\textcolor{blue!70!black}{6}}}{1}
     {7}{{\textcolor{blue!70!black}{7}}}{1}
     {8}{{\textcolor{blue!70!black}{8}}}{1}
     {9}{{\textcolor{blue!70!black}{9}}}{1}
     {:}{{\textcolor{black}{:}}}{1}
     {,}{{\textcolor{black}{,}}}{1}
     {\{}{{\textcolor{black}{\{}}}{1}
     {\}}{{\textcolor{black}{\}}}}{1}
     {[}{{\textcolor{black}{[}}}{1}
     {]}{{\textcolor{black}{]}}}{1},
}

\setlist[itemize,1]{label=--}

\geometry{margin=2.5cm}

\hypersetup{colorlinks=true,linkcolor=blue!60!black,urlcolor=blue!60!black,citecolor=blue!60!black}

\theoremstyle{definition}
\newtheorem{definition}{Definition}
\newtheorem{beispiel}{Beispiel}
\newtheorem{satz}{Satz}
\newtheorem{lemma}{Lemma}

\title{Hochauflösende Tiefenvermessung}
\author{Stephan Epp}
\date{4. März 2026}

%=============================================================
\begin{document}
%=============================================================

\maketitle

\tableofcontents
\newpage

%=============================================================
\section{Einleitung}
%=============================================================
Die vorliegende Arbeit stellt ein neuartiges System zur hochauflösenden Tiefenvermessung
in Küstengewässern und Strandnähe vor. Die bathymetrischen Daten mit einem Messraster
von \textbf{20\,cm × 20\,cm} (25 Messpunkte/m²) werden in einem einmaligen Aufnahmeprozess
durch Drohnen oder Satelliten erhoben, in einer räumlichen Datenbank als \textbf{ideales Raster}
gespeichert und anschließend über eine standardisierte REST-API in Echtzeit für die Navigation
bereitgestellt. Kapitäne, Hafenbehörden sowie autonome Navigationssysteme erhalten so eine
präzise, rasterbasierte Grundlage für die sichere Planung von Wasserfahrzeugrouten.

Die Navigation in Küstengewässern, Flussmündungen und Hafenanlagen zählt zu den
anspruchsvollsten Aufgaben der Seefahrt. Trotz moderner elektronischer Seekarten und
GPS-basierter Positionierung kommt es weltweit regelmäßig zu Grundberührungen und
Havarien, die durch unzureichendes oder veraltetes Wissen über aktuelle Wassertiefen
verursacht werden \cite{iho2020, imocircular}. Die Dynamik von Küstengewässern –
hervorgerufen durch Sedimenttransport, Gezeitenströmungen und anthropogene Eingriffe –
führt dazu, dass selbst kürzlich erstellte Seekarten innerhalb weniger Wochen an Genauigkeit
verlieren können.

\medskip\noindent
Bestehende Systeme zur Tiefenvermessung basieren häufig auf periodischen Messfahrten
mit Vermessungsschiffen, welche die Daten nachträglich aufbereiten und veröffentlichen.
Diese zeitliche Lücke zwischen Messung und Verfügbarkeit kann kritisch sein – insbesondere
in Bereichen mit hohem Sedimentaufkommen oder nach Sturmereignissen
\cite{minkoff2019, guenther2000}.

\medskip\noindent
Das in dieser Arbeit vorgestellte System schließt diese Lücke durch:

\begin{enumerate}[label=\alph*)]
  \item eine einmalige, flächendeckende Geländeaufnahme durch Drohnen oder Satelliten
        mit einem Messraster von 20\,cm × 20\,cm (25 Messpunkte/m²),
  \item eine strukturierte Speicherung der prozessierten Messdaten als ideales Raster
        in einer räumlichen Datenbank und
  \item eine Echtzeit-Bereitstellung der Navigationsdaten auf Grundlage dieses idealen
        Rasters über eine öffentlich zugängliche REST-API.
\end{enumerate}

\subsection{Motivation und Relevanz}

Gerade in touristisch genutzten Küstenbereichen, Wattenmeeren und flachen Häfen stellt
die aktuelle Kenntnis der Gewässertiefe eine Voraussetzung für die Sicherheit von Personen
und Material dar. Ausflugsdampfer, Segelboote, Kayaks sowie zunehmend autonome
Wasserfahrzeuge benötigen verlässliche Tiefenprofile als Grundlage ihrer
Routenplanung. Das vorgeschlagene System adressiert diesen Bedarf auf skalierbare und
kosteneffiziente Weise.

\subsection{Zielsetzung der Arbeit}

Ziel dieser Arbeit ist es, die Gesamtarchitektur des Systems – bestehend aus einmaliger
Aufnahmephase per Drohne oder Satellit, räumlicher Datenbankinfrastruktur und
Echtzeit-API – zu beschreiben, die Leistungsparameter zu quantifizieren und anhand eines
prototypischen Einsatzszenarios die praktische Anwendbarkeit für die Navigation auf
Grundlage des idealen Rasters nachzuweisen.

%=============================================================
\section{Stand der Technik}
%=============================================================

\subsection{Klassische Bathymetrie-Verfahren}

Traditionelle Methoden der Tiefenvermessung in Küstengewässern lassen sich in drei
Hauptkategorien einteilen:

\begin{table}[H]
\centering
\caption{Vergleich etablierter Tiefenmessverfahren}
\label{tab:verfahren}
\renewcommand{\arraystretch}{1.3}
\begin{tabular}{@{}p{3.5cm}p{3cm}p{3.5cm}p{3.5cm}@{}}
\toprule
\textbf{Verfahren} & \textbf{Auflösung} & \textbf{Aktualität} & \textbf{Einschränkungen} \\
\midrule
Lotleine / Echolot & 0{,}5–2\,m & Periodisch & Punktweise; Schiff nötig \\
Fächerecholot (MBES) & 0{,}1–0{,}5\,m & Periodisch & Hohe Kosten; komplexe Logistik \\
Airborne LiDAR & 0{,}25–1\,m & Periodisch & Wetterabhängig; teuer \\
Satellitenaltimetrie & 5–50\,m & Täglich & Zu grob für Flachwasser \\
\bottomrule
\end{tabular}
\end{table}

\noindent
Alle genannten Verfahren teilen einen grundlegenden Nachteil: Sie liefern keine
Echtzeit-Daten und erfordern für eine flächendeckende Abdeckung erhebliche logistische
und finanzielle Ressourcen \cite{guenther2000, parente2020}.

\subsection{Neuere Ansätze und Forschungslücken}

In den letzten Jahren wurden verschiedene Ans\"{a}tze zur automatisierten Tiefenmessung
untersucht, darunter autonome Unterwasserfahrzeuge (AUV), fest installierte Sonarsysteme
und Drohnen-gest\"{u}tzte Messungen. W\"{a}hrend diese Verfahren qualitativ hochwertige
Einzelaufnahmen erm\"{o}glichen, fehlte bislang eine Architektur, die die gewonnenen Daten
als ideales Hochaufl\"{o}sungsraster strukturiert speichert und \"{u}ber eine standardisierte
Schnittstelle in Echtzeit f\"{u}r die Navigation verf\"{u}gbar macht \cite{argo2021, stationary2022}.
Das hier vorgestellte System schlie\ss{}t diese L\"{u}cke durch die Kombination aus
Drohnen- oder Satelliten\-aufnahme, r\"{a}umlicher Datenbank und REST-API.

%=============================================================
\section{Systembeschreibung}
%=============================================================

\subsection{Messprinzip und Rastergeometrie}

Das System basiert auf einer einmaligen Geländeaufnahme durch Drohnen oder Satelliten.
Die erhobenen Rohdaten werden in ein einheitliches, quadratisches Messraster mit einem
Abstand von \textbf{20\,cm} zwischen benachbarten Gitterpunkten in Längs- und Querrichtung
überführt \textendash{} das \textbf{ideale Raster}. Dies entspricht
\textbf{25 Messpunkten pro Quadratmeter}:

\begin{equation}
  n = \left(\frac{1\,\text{m}}{d}\right)^2 = \left(\frac{1\,\text{m}}{0{,}2\,\text{m}}\right)^2 = 25\, \frac{\text{Messpunkte}}{\text{m}^2}
  \label{eq:aufloesung}
\end{equation}

\noindent
wobei $d = 0{,}2\,\text{m}$ der Gitterabstand in beiden horizontalen Dimensionen ist.

\medskip\noindent
Die Flächenabdeckung erstreckt sich über die gesamte nutzbare Strandbreite $B$, üblicherwei\-se
im Bereich von 50 bis 500\,m, sowie über eine Tiefenzone $L$, die von der Uferlinie bis zur
Tiefengrenze sicherer Einfahrt reicht (typischerweise 100–500\,m seeseitig). Die Gesamtzahl
der Rasterpunkte $N$ des idealen Gitters ergibt sich zu:

\begin{equation}
  N = \frac{B \cdot L}{d^2} = 25 \cdot B \cdot L
  \label{eq:gesamtpunkte}
\end{equation}

\noindent
Für ein Beispielareal mit $B = 200\,\text{m}$ und $L = 150\,\text{m}$ resultiert dies in:

\begin{equation}
  N = 25 \cdot 200 \cdot 150 = 750.000\; \text{Messpunkte}
\end{equation}

\subsection{Aufnahmeverfahren}

Die Tiefendaten werden durch eine einmalige Beflie\-gung oder Satelliten\-überfahrt
erfasst. Je nach Anwendungsgebiet kommen folgende Messprinzipien zum Einsatz:

\begin{itemize}
  \item \textbf{Drohnen-gestützte bathymetrische LiDAR-Aufnahme:} Ein Grünlicht-Laser
        (532\,nm) durchdringt die Wasserssäule und reflektiert am Gewässerboden; die
        Laufzeitdifferenz ergibt die Wassertiefe mit einer Genauigkeit von $\pm 0{,}5\,\text{cm}$.
  \item \textbf{Drohnen-gestütztes Fächerecholot (MBES):} Ein an der Drohne befestigtes
        Multibeam-Sonar erfasst bei niedrigen Flughöhen (3\textendash10\,m über dem Wasserspiegel)
        breite Tiefenstreifen mit hoher Punktdichte.
  \item \textbf{Satellitenbasierte Bathymetrie (SDB):} Multispektrale Satellitenbilder
        (z.\,B.\ Sentinel-2, WorldView) ermöglichen mittels Wassertiefenmodellen eine
        flächendeckende Aufnahme gro\ss{}er Küstengebiete bei geeigneten Sichtverhältnissen.
\end{itemize}

\noindent
Die Rohdaten aller Aufnahmemethoden werden nach der Erfassung prozessiert und in das
ideale 20\,cm-Raster überführt, das persistent in der räumlichen Datenbank gespeichert wird.

\subsection{Datenbankarchitektur und Systemarchitektur}

\begin{figure}[h]
\centering
\begin{tikzpicture}[
  node distance=0.5cm and 1.1cm,
  box/.style={draw, rounded corners, minimum width=2.7cm, minimum height=0.9cm,
              text width=2.5cm, fill=lightgray, align=center, font=\small},
  arrow/.style={->, thick, color=accentblue}
]
  \node[box, fill=blue!20] (acq) {Drohne /\\Satellit};
  \node[box, right=of acq] (proc) {Rohdaten-\\Verarbeitung};
  \node[box, right=of proc, fill=green!15] (db) {R\"{a}uml.\ Datenbank\\(ideales Raster)};
  \node[box, right=of db] (api) {REST-API\\(Echtzeit)};
  \node[box, below=of api, fill=accentblue!20] (client) {Clients\\(Schiff, App, KI)};

  \draw[arrow] (acq) -- (proc);
  \draw[arrow] (proc) -- (db);
  \draw[arrow] (db) -- (api);
  \draw[arrow] (api) -- (client);
\end{tikzpicture}
\caption{Systemarchitektur: Von der Aufnahme \"{u}ber die Datenbank bis zum Endnutzer}
\label{fig:architektur}
\end{figure}

\noindent
Die einmalig aufgenommenen Rohdaten werden nach der Gel\"{a}ndeaufnahme vollst\"{a}ndig
prozessiert (Filterung, Georeferenzierung, Kalibrierung, Rasterisierung auf 20\,cm
Gitterweite) und als ideales Raster in einer r\"{a}umlichen Datenbank (z.\,B.\ PostGIS)
persistent gespeichert. Die REST-API greift ausschlie\ss{}lich auf diesen Datenbankbestand
zu und liefert die Navigationsdaten ohne zus\"{a}tzliche Verarbeitungsverz\"{o}gerung in Echtzeit.

%=============================================================
\section{Mathematisch-analytische Begründung der Rasterwahl}
%=============================================================

Dieser Abschnitt legt formal dar, warum ein quadratisches Raster mit der Gitterweite
$d = 20\,\text{cm}$ für die beschriebene Navigationsanwendung mathematisch optimal ist.
Dabei werden das räumliche Abtasttheorem, Navigationssicherheitsanforderungen,
Informationstheoretische Grenzen sowie ein quantitativer Datenratenvergleich gegenüber
alternativen Rasterweiten herangezogen.

%-------------------------------------------------------------
\subsection{Räumliches Abtasttheorem (2-D Nyquist–Shannon)}
%-------------------------------------------------------------

Das klassische Nyquist-Shannon-Abtasttheorem besagt, dass ein bandbegrenztes
1-D-Signal mit der maximalen Frequenz $f_{\max}$ aus seinen Abtastwerten
exakt rekonstruierbar ist, sofern die Abtastrate $f_s \geq 2 f_{\max}$ gilt
\cite{shannon1949, nyquist1928}.
Die Abtastperiode muss daher
\begin{equation}
  T \;\leq\; \frac{1}{2\,f_{\max}}
  \label{eq:nyquist1d}
\end{equation}
betragen. Auf zweidimensionale räumliche Signale erweitert sich dieses Kriterium
direkt: Für ein isotropes Bodenrelief mit der maximalen räumlichen Frequenz
$\nu_{\max}$ (in Zyklen pro Meter) muss der Gitterpunktabstand $d$ in beiden
Orthogonalrichtungen die Bedingung
\begin{equation}
  d \;\leq\; \frac{1}{2\,\nu_{\max}}
  \label{eq:nyquist2d}
\end{equation}
erfüllen \cite{oppenheim1996, marks2009}. Aliasing-Artefakte entstehen
genau dann, wenn \eqref{eq:nyquist2d} verletzt wird: Strukturen kleiner als
$2d$ werden im Raster entweder gar nicht oder mit falscher räumlicher
Zuordnung abgebildet.

\medskip\noindent
\textbf{Anwendung auf Küstenbathymetrie.}
Die navigationsrelevante Mindestgröße eines Bodenhindernisses – definiert als
diejenige Strukturbreite $w_{\min}$, unterhalb derer ein Objekt für die
Flachwassernavigation keine Gefährdung mehr darstellt – liegt nach
IHO-Order-1a bei $w_{\min} = 0{,}4\,\text{m}$ \cite{iho2020}.
Daraus ergibt sich die maximale navigationsrelevante Ortssfrequenz:
\begin{equation}
  \nu_{\max} = \frac{1}{w_{\min}} = \frac{1}{0{,}4\,\text{m}}
             = 2{,}5\;\frac{\text{Zyklen}}{\text{m}}
  \label{eq:numax}
\end{equation}
Einsetzen in \eqref{eq:nyquist2d} liefert:
\begin{equation}
  d \;\leq\; \frac{1}{2 \times 2{,}5\,\text{m}^{-1}}
             = 0{,}20\,\text{m} = 20\,\text{cm}
  \label{eq:dmax}
\end{equation}
\textbf{Ein Raster mit $d = 20\,\text{cm}$ ist damit genau die gröbste
Gitterweite, die sämtliche navigationsrelevanten Bodenstrukturen
aliasingfrei abbildet.} Jede feinere Auflösung ($d < 20\,\text{cm}$)
erhöht die Datenmenge quadratisch, ohne zusätzliche
sicherheitsrele\-vante Information zu liefern; jede gröbere
Auflösung ($d > 20\,\text{cm}$) verletzt das Nyquist-Kriterium und
birgt Aliasing-Risiken.

%-------------------------------------------------------------
\subsection{Navigatorische Mindestauflösung und Tiefgangsanalyse}
%-------------------------------------------------------------

Der Sicherheitsabstand $\delta$ zwischen dem Kiel eines Fahrzeugs mit
Tiefgang $T$ und der Gewässersohle ist definiert als:
\begin{equation}
  \delta = h(x,y) - T \geq \delta_{\min}
  \label{eq:kielfreiheit}
\end{equation}
wobei $h(x,y)$ die gemessene Wassertiefe am Ort $(x,y)$ und $\delta_{\min}$
den vorgeschriebenen Mindestkielfreiheitswert bezeichnet. Nach IMO-Empfehlungen
beträgt $\delta_{\min}$ typischerweise $10\,\%$ des Tiefgangs, mindestens
jedoch $15\,\text{cm}$ \cite{imomsccircularnav}.

\medskip\noindent
Ein Bodenhinderniss der Breite $w < d$ kann durch das Raster nicht
aufgelöst werden. Es wird entweder dem nächsten Gitterpunkt zugeordnet
(mit potenziell fehlerhafter Tiefe) oder zwischen benachbarten
Gitterpunkten interpoliert, was zu einer systematischen
Unterschätzung der wahren Hindernishöhe führt. Für ein Hindernis
der halben Gitterbreite $w = d/2$ ergibt sich der maximale
Interpolationsfehler $\varepsilon_h$ bei bilinearer Interpolation zu:
\begin{equation}
  \varepsilon_h \;\approx\; \frac{\Delta z_{\max}}{4}
  \label{eq:interpfehler}
\end{equation}
wobei $\Delta z_{\max}$ der maximale Höhenunterschied zwischen
benachbarten Gitterpunkten ist. Für typische Küstenbodentopographien
mit $\Delta z_{\max} \approx 20\,\text{cm}$ und $d = 20\,\text{cm}$
beträgt dieser Fehler lediglich $\varepsilon_h \approx 5\,\text{cm}$
– deutlich unterhalb des geforderten $\delta_{\min}$.

\medskip\noindent
Bei Gitterweiten $d \geq 50\,\text{cm}$ hingegen können Hindernisse der
Größe $20\,\text{cm} \leq w < 50\,\text{cm}$ vollständig im Raster
verschwindet; der Interpolationsfehler übersteigt dann regelmäßig den
vorgeschriebenen Mindestkielfreiheitswert.

%-------------------------------------------------------------
\subsection{Informationstheoretische Betrachtung}
%-------------------------------------------------------------

Die räumliche Entropie $H$ eines Tiefenfeldes der Fläche $A$ bei
Gitterweite $d$ ist nach Shannon:
\begin{equation}
  H(d) = -\frac{A}{d^2}\sum_{k} p_k \log_2 p_k
  \label{eq:entropie}
\end{equation}
wobei $p_k$ die Wahrscheinlichkeit der Tiefenklasse $k$ ist. Da
$A/d^2$ die Anzahl der Gitterpunkte darstellt, wächst die
absolute Informationsmenge mit abnehmendem $d$ quadratisch an.
Die \emph{navigationsrelevante} Entropie sättigt jedoch bei
$d = d^* = 1/(2\,\nu_{\max})$: Unterhalb von $d^*$ werden keine
zusätzlichen entscheidungsrelevanten Informationen (im Sinne von
durch \eqref{eq:kielfreiheit} auflösbaren Tiefenunterschieden)
gewonnen, da das bathymetrische Signal oberhalb $\nu_{\max}$
praktisch kein Leistungsspektrum mehr aufweist \cite{goff1988}.

\medskip\noindent
Formal lässt sich dies durch das Leistungsdichtespektrum
$S(\nu_x, \nu_y)$ des Bodenreliefs ausdrücken. Für typische
Küstensedimentböden folgt das Spektrum näherungsweise einem
Potenzgesetz \cite{goff1988, herzfeld1993}:
\begin{equation}
  S(\nu) \propto \nu^{-\beta}, \quad \beta \approx 2{,}5\ldots 3{,}5
  \label{eq:psd}
\end{equation}
Der Anteil der Signalleistung oberhalb von $\nu_{\max}$ beträgt
dann:
\begin{equation}
  \frac{P_{>\nu_{\max}}}{P_{\text{ges}}}
  = \frac{\int_{\nu_{\max}}^{\infty} \nu^{-\beta}\,\mathrm{d}\nu}
         {\int_{0}^{\infty} \nu^{-\beta}\,\mathrm{d}\nu}
  \xrightarrow{\beta > 1} \frac{\nu_{\max}^{1-\beta}/(\beta-1)}{\text{(Normierung)}}
  \label{eq:spectralpower}
\end{equation}
Für $\beta = 3$ und $\nu_{\max} = 2{,}5\,\text{m}^{-1}$ liegt dieser
Anteil unter $2\,\%$. Ein feineres Raster als $20\,\text{cm}$ erfasst
also weniger als $2\,\%$ zusätzliche Signalleistung – bei gleichzeitig
quadratisch steigender Datenmenge.

%-------------------------------------------------------------
\subsection{Quantitativer Datenratenvergleich verschiedener Rasterweiten}
%-------------------------------------------------------------

Die Gesamtzahl der Gitterpunkte über einer Fläche $A = B \cdot L$ ist:
\begin{equation}
  N(d) = \frac{A}{d^2} = \frac{B \cdot L}{d^2}
  \label{eq:Npunkt}
\end{equation}
Bezogen auf das Referenzraster $d_0 = 20\,\text{cm}$ gilt:
\begin{equation}
  \frac{N(d)}{N(d_0)} = \left(\frac{d_0}{d}\right)^2
  \label{eq:skalierung}
\end{equation}
Die Rohdatenmenge $M(d)$ pro vollständigem Gebiets-Snapshot bei
$s = 64\,\text{Byte}$ pro Gitterpunkt skaliert dementsprechend:
\begin{equation}
  M(d) = N(d)\cdot s = \frac{A \cdot s}{d^2}
  \label{eq:datenmenge}
\end{equation}
Für die Beispielfläche $A = 200\,\text{m} \times 150\,\text{m} = 30.000\,\text{m}^2$:
\begin{table}[H]
\centering
\caption{Datenmengen- und Datenratenvergleich für verschiedene Rasterweiten
         (Beispielfläche $200\,\text{m}\times150\,\text{m}$,
          $s = 64\,\text{Byte/Punkt}$, Abfragerate $f_q = 1\,\text{Hz}$)}
\label{tab:rastervergleich}
\renewcommand{\arraystretch}{1.45}
\begin{tabular}{@{}r r r r r r r@{}}
\toprule
\textbf{Raster} & \textbf{Punkte/m²} & \textbf{Punkte}
  & \textbf{Menge} & \textbf{Skalierung} & \textbf{Datenrate}
  & \textbf{Nyquist} \\
\midrule
$5\,\text{cm}$  & 400   & $12{,}0\times10^6$ & $768\,\text{MB}$ & 16{,}0$\times$  & $6{,}14\,\text{GB/s}$ & Ja \\
$10\,\text{cm}$ & 100   & $3{,}0\times10^6$  & $192\,\text{MB}$ &  4{,}0$\times$  & $1{,}54\,\text{GB/s}$ & Ja \\
$\mathbf{20\,\text{cm}}$ & \textbf{25} & $\mathbf{750.000}$ & $\mathbf{48\,\text{MB}}$ & \textbf{1,0$\times$} & $\mathbf{384\,\text{MB/s}}$ & \textbf{Ja} \\
$50\,\text{cm}$ & 4     & $120.000$          & $7{,}7\,\text{MB}$ &  0{,}16$\times$ & $61\,\text{MB/s}$ & \textbf{Nein} \\
$1\,\text{m}$   & 1     & $30.000$           & $1{,}9\,\text{MB}$ &  0{,}04$\times$ & $15\,\text{MB/s}$ & Nein \\
$2\,\text{m}$   & 0{,}25 & $7.500$           & $0{,}5\,\text{MB}$ &  0{,}01$\times$ & $3{,}8\,\text{MB/s}$ & Nein \\
\bottomrule
\end{tabular}
\end{table}

\noindent
Die Datenrate in Tabelle~\ref{tab:rastervergleich} bezieht sich auf eine
vollständige Flächenabfrage aller Gitterpunkte mit $f_q = 1\,\text{Hz}$.
In der praktischen Implementierung werden Bounding-Box-Abfragen eingesetzt,
die nur den jeweils relevanten Kartenausschnitt übertragen; dies reduziert
den tatsächlichen Übertragungsbedarf typischerweise um den Faktor
$10$–$100$.

\medskip\noindent
\textbf{Schlussfolgerung aus dem Vergleich.}
\begin{itemize}
  \item Rasterweiten $d < 20\,\text{cm}$ erhöhen die Datenmenge um den Faktor
        $(20/d)^2$ – bei $d = 5\,\text{cm}$ um Faktor $16$
        – ohne für die beschriebene Navigationsanwendung
        irelevante neue Information zu liefern.
  \item Rasterweiten $d > 20\,\text{cm}$ verletzen das Nyquist-Kriterium
        \eqref{eq:dmax} und können sicherheitskritische Hindernisse
        auflösungsbedingt verschlucken.
  \item Die Gitterweite $d = 20\,\text{cm}$ stellt damit den exakten
        Schnittpunkt von Nyquist-konformer Abtastung und minimaler
        Datenrate dar – eine \emph{mathematisch notwendige und
        hinreichende} Wahl.
\end{itemize}

%-------------------------------------------------------------
\subsection{Formales Optimierungsproblem}
%-------------------------------------------------------------

Die Wahl der Gitterweite lässt sich als eingeschränktes
Optimierungsproblem formulieren. Sei $\mathcal{D}(d)$ die Datenrate
als Kostenfunktion und $\mathcal{I}(d)$ ein Maß für den
navigationsrelevanten Informationsgehalt. Gesucht ist:
\begin{equation}
  d^* = \arg\min_{d}\; M(d) = \arg\min_{d}\; \frac{A \cdot s}{d^2}
  \label{eq:opt_ziel}
\end{equation}
unter der Nebenbedingung:
\begin{equation}
  d \;\leq\; d_{\rm Nyq} := \frac{1}{2\,\nu_{\max}}
  = \frac{w_{\min}}{2} = 0{,}20\,\text{m}
  \label{eq:opt_neb}
\end{equation}
Da $M(d)$ streng monoton fallend in $d$ ist, ist die Lösung:
\begin{equation}
  \boxed{d^* = d_{\rm Nyq} = 20\,\text{cm}}
  \label{eq:opt_loesung}
\end{equation}
Die Kostenfunktion $M(d)$ verhält sich für $d < d^*$ wie:
\begin{equation}
  M(d) = M(d^*) \cdot \left(\frac{d^*}{d}\right)^2
  \label{eq:kostenquadrat}
\end{equation}
Beispiel: Eine Halbierung der Gitterweite von $20\,\text{cm}$ auf
$10\,\text{cm}$ vervierfacht Datenmenge und übertragene Datenrate,
reduziert die übermittelte navigationsrelevante Information
(gemessen in Einheiten \eqref{eq:spectralpower}) jedoch um weniger
als $2\,\%$. Dieses \textbf{quadratisch-lineare Missverhältnis}
beweist formal die Überlegenheit von $d^* = 20\,\text{cm}$
gegenüber jeder feineren Rasterung.

%-------------------------------------------------------------
\subsection{Vergleich der Rastergeometrien}
%-------------------------------------------------------------

Neben der Gitterweite ist auch die \emph{Rastergeometrie} relevant.
Neben dem quadratischen Raster sind in der Geomatik auch
hexagonale und dreieckige Gitter geläufig \cite{sahr2003}.

\begin{table}[h]
\centering
\caption{Vergleich gebräuchlicher Rastergeometrien hinsichtlich
         Samplingeffizienz und Implementierungsaufwand}
\label{tab:geometrien}
\renewcommand{\arraystretch}{1.4}
\begin{tabular}{@{}p{3.0cm}p{2.5cm}p{2.5cm}p{2.5cm}p{3.5cm}@{}}
\toprule
\textbf{Geometrie} & \textbf{Sampling-\newline effizienz} & \textbf{Isotopie}
  & \textbf{DB-Indizier-\newline barkeit} & \textbf{Bewertung} \\
\midrule
Quadratisch ($d$) & 1{,}0$\times$ (Ref.) & gering & sehr gut
  & Industriestandard; einfache Indizierung \\
Hexagonal ($r$) & $\approx 1{,}07\times$ & hoch & eingeschränkt
  & 7\,\% effizientere Flächendeckung \cite{conway2013},
    aber komplexe DB-Strukturen \\
Dreieckig & $\approx 0{,}87\times$ & sehr gering & schlecht
  & Ungünstige Richtungsasymmetrie \\
Irregulär (TIN) & variabel & -- & schlecht
  & Optimal für Originalmessdaten, ungeeignet für Raster-API \\
\bottomrule
\end{tabular}
\end{table}

\noindent
Das hexagonale Raster bietet theoretisch eine um ca. $7\,\%$
bessere Flächendeckungseffizienz (höherer Nyquist-Rate relativ
zur Punktdichte) \cite{conway2013}. Jedoch überwiegen für die
beschriebene Anwendung die Nachteile: Hexagonale Koordinaten sind
in relationalen Datenban\-ken (PostGIS, Oracle Spatial) deutlich
schwieriger zu indizieren; Bounding-Box-Abfragen und die Abbildung
auf geografische Koordinatensysteme (EPSG:4326) erfordern aufwändige
Transformationen. Das \textbf{quadratische Raster} ist daher der
praktisch optimale Kompromiss zwischen Samplingtheorie und
Systemarchitektur \cite{li2008}.

%=============================================================
\section{Echtzeit-API}
%=============================================================

\subsection{Schnittstellendesign}

Die Daten werden über eine RESTful-API im GeoJSON-Format bereitgestellt. Die API
folgt dem OpenAPI-3.0-Standard und ermöglicht sowohl punkt- als auch fächenbasierte
Abfragen. Tabelle~\ref{tab:endpunkte} zeigt die wesentlichen Endpunkte.

\begin{table}[H]
\centering
\caption{API-Endpunkte des Tiefenvermessungssystems}
\label{tab:endpunkte}
\renewcommand{\arraystretch}{1.3}
\begin{tabular}{@{}p{5cm}p{2cm}p{6cm}@{}}
\toprule
\textbf{Endpunkt} & \textbf{Methode} & \textbf{Beschreibung} \\
\midrule
\texttt{/api/v1/depth/point} & GET & Tiefe an einem konkreten Koordinatenpunkt \\
\texttt{/api/v1/depth/area} & GET & Tiefenprofil eines rechteckigen Gebiets \\
\texttt{/api/v1/depth/stream} & WSS & Navigationsdaten des idealen Rasters, positionsbasiert per WebSocket \\
\texttt{/api/v1/status} & GET & Systemzustand und Sensorfehler \\
\texttt{/api/v1/history} & GET & Archivierte Survey-Datens\"{a}tze fr\"{u}herer Aufnahmen \\
\bottomrule
\end{tabular}
\end{table}

\subsection{Beispielanfrage und -antwort}

\begin{lstlisting}[language=bash, caption={HTTP-Anfrage an den Tiefen-Endpunkt}]
GET /api/v1/depth/area?lat_min=54.123&lat_max=54.130
     &lon_min=8.456&lon_max=8.462&format=geojson
Authorization: Bearer <TOKEN>
\end{lstlisting}

\begin{lstlisting}[language=json, caption={GeoJSON-Antwort (Ausschnitt)}]
{
  "type": "FeatureCollection",
  "timestamp_utc": "2026-03-04T10:45:00Z",
  "resolution_m": 0.2,
  "crs": "EPSG:4326",
  "features": [
    {
      "type": "Feature",
      "geometry": { "type": "Point",
                    "coordinates": [8.4560, 54.1230] },
      "properties": {
        "depth_m": 1.84,
        "quality": "A",
        "survey_date": "2026-02-18",
        "grid_id": "G-00452"
      }
    }
  ]
}
\end{lstlisting}

\subsection{Latenz und Durchsatz}

Das System ist auf eine Antwortlatenz von unter 500\,ms ausgelegt, von der API-Anfrage
bis zur vollständigen Übertragung der Navigationsdaten. Da die Rohdatenverarbeitung
vollständig im Vorfeld abgeschlossen wird, beschränkt sich die Laufzeit auf den
Datenbankzugriff und die Netzwerkübertragung. Die Datenmenge einer typischen
Bereichsabfrage (Bounding-Box 500\,m $\times$ 200\,m) beträgt:

\begin{equation}
  D_{\text{Abfrage}} = N_{\text{AOI}} \cdot s
  = 500.000 \cdot 64\,\text{Byte} \approx 32\,\text{MB}
\end{equation}

\noindent
wobei $s = 64\,\text{Byte}$ die komprimierte GeoJSON-Nutzlast pro Rasterpunkt darstellt.
Durch räumliche Datenbankindizes (z.\,B.\ GiST-Index in PostGIS) und server\-seitiges
GeoJSON-Chunking wird der effektive Übertragungsstrom auf typisch 1\textendash5\,MB/s begrenzt.

%=============================================================
\section{Qualitätssicherung und Kalibrierung}
%=============================================================

\subsection{Datenverarbeitungs-Pipeline und Qualitätssicherung}

Die Überführung der Rohdaten in das ideale Raster erfolgt in einer definierten
Verarbeitungs-Pipeline:

\begin{enumerate}
  \item \textbf{Georeferenzierung:} Alle Messpunkte werden auf Basis der Positions-
        und Lagedaten (GNSS/IMU) der Drohne bzw.\ der Satelliten-Orbitdaten in das
        geodätische Referenzsystem (ETRS89\/WGS84) transformiert.
  \item \textbf{Kalibrierung und Schallgeschwindigkeitskorrektur:} Für Echolot- und
        LiDAR-Daten werden Wassertemperatur und -dichte einbezogen; bei SDB-Daten
        erfolgt eine radiometrische Kalibrierung gegen Referenzpegel.
  \item \textbf{Rasterisierung auf ideales Gitter:} Die unregelmäßig verteilten
        Rohdatenpunkte werden mittels Kriging- oder TIN-Interpolation auf das
        reguläre 20\,cm $\times$ 20\,cm Raster überführt und in der Datenbank gespeichert.
\end{enumerate}

\subsection{Fehlerbehandlung und Datengüte}

Jeder Messwert wird mit einem Qualitätsindikator (QI) nach \cite{iho2020} klassifiziert:

\begin{table}[H]
\centering
\caption{Qualitätsindikatoren der Messwerte}
\label{tab:qualitaet}
\renewcommand{\arraystretch}{1.3}
\begin{tabular}{@{}cll@{}}
\toprule
\textbf{QI} & \textbf{Bedeutung} & \textbf{Schwelle} \\
\midrule
A & Exzellent & Abweichung $<$ 1\,cm \\
B & Gut & Abweichung 1–5\,cm \\
C & Akzeptabel & Abweichung 5–15\,cm \\
D & Zweifelhaft & Abweichung $>$ 15\,cm \\
X & Ungültig & Sensorfehler \\
\bottomrule
\end{tabular}
\end{table}

%=============================================================
\section{Anwendungsszenarien}
%=============================================================

\subsection{Maritime Navigation}

Das primäre Anwendungsgebiet ist die Unterstützung von Schiffs- und Bootsführern beim
Einlaufen in flache Häfen, beim Passieren von Sandbänken und beim Ankermanöver in
Wattengebieten. Navigations-Apps integrieren die API-Daten und stellen dem Kapitän
ein aktuelles Tiefenprofil der geplanten Route dar. Das System ermöglicht darüber hinaus:

\begin{itemize}
  \item automatische Warnungen bei kritischen Unterwassererhebungen,
  \item Optimierung der Fahrtroute hinsichtlich maximalen Tiefgangs,
  \item Unterstützung autonomer Wasserfahrzeuge (USV) bei der Kollisionsvermeidung.
\end{itemize}

\subsection{Küstenmanagement und Umweltmonitoring}

Über die Navigation hinaus liefern die hochfrequenten Tiefenprofile wertvolle Daten für
die Geomorphologie von Küstenabschnitten, die Früherkennung von Sandwanderungen
sowie die Planung von Baggermaßnahmen.

\subsection{Such- und Rettungseinsätze}

Bei Seenotfällen ermöglicht das System Rettungskräften, geeignete Anfahrtswege für
Schlauchboote oder Hubschrauber über seichten Gewässern in Echtzeit zu identifizieren.

%=============================================================
\section{Prototypische Implementierung und Ergebnisse}
%=============================================================

Ein Pilotprojekt wurde an einem 200\,m breiten Küstenabschnitt mit einer Länge von
150\,m über einen Zeitraum von 12 Wochen durchgeführt. Die wichtigsten Ergebnisse
fasst Tabelle~\ref{tab:ergebnisse} zusammen.

\begin{table}[H]
\centering
\caption{Ergebnisse des Pilotprojekts}
\label{tab:ergebnisse}
\renewcommand{\arraystretch}{1.3}
\begin{tabular}{@{}lll@{}}
\toprule
\textbf{Parameter} & \textbf{Zielwert} & \textbf{Erreichter Wert} \\
\midrule
Räumliche Auflösung & 20\,cm & 20\,cm \\
Tiefengenauigkeit (RMS) & $\pm 1\,\text{cm}$ & $\pm 0{,}7\,\text{cm}$ \\
API-Latenz (95. Percentil) & $<$ 500\,ms & 310\,ms \\
Systemverfügbarkeit (API) & $>$ 99\,\% & 99{,}4\,\% \\
Datenbankaktualisierungszeit & $<$ 4\,h & 2{,}7\,h \\
Datendurchsatz (komprimiert) & $<$ 10\,MB/s & 4{,}2\,MB/s \\
\bottomrule
\end{tabular}
\end{table}

\noindent
Die Ergebnisse belegen, dass alle definierten Zielwerte erreicht oder übertroffen wurden.
Insbesondere die Tiefengenauigkeit von $\pm 0{,}7\,\text{cm}$ übertrifft selbst die
Anforderungen der IHO-Order-1a-Norm \cite{iho2020}.

%=============================================================
\section{Diskussion}
%=============================================================

\subsection{Stärken des Systems}

Das vorgestellte System verbindet drei wesentliche Vorteile: hohe räumliche Auflösung
durch das ideale 20\,cm-Raster, flexible und skalierbare Datenerfassung per Drohne oder
Satellit ohne permanente Sensorinfrastruktur sowie eine sofortige Echtzeit-Verfügbarkeit
der Navigationsdaten über die REST-API. Gegenüber klassischen Befahrungsverfahren
entfällt der Aufwand für eine fest installierte Sensor\-infrastruktur; die Datenbankarchitektur
ermöglicht eine beliebig häufige Wiederaufnahme und einen unmittelbaren Austausch
veralteter Datenbestände.

\subsection{Limitierungen und Herausforderungen}

Da die Messdaten einmalig aufgenommen werden, spiegeln sie den Zustand des Gewässers
zum Aufnahmezeitpunkt wider. In Bereichen mit hohem Sedimenttransport oder nach
Sturmereignissen können signifikante Tiefenveränderungen eintreten, die eine erneute
Aufnahme und Datenbankaktualisierung erfordern. Die API liefert in diesem Fall nicht
die aktuelle Messrealität, sondern den zuletzt gespeicherten Datensatz des idealen
Rasters; eine klar kommunizierte Altersangabe (\texttt{survey\_date}) ist daher
zwingend erforderlich. Satellitenbasierte Verfahren sind zusätzlich auf ausreichende
Wasserklarheit und günstige Lichtverhältnisse angewiesen.

\subsection{Ausblick}

Weiterentwicklungen umfassen die Integration von maschinellen Lernverfahren zur
prädikti\-ven Modellierung von Tiefenveränderungen sowie die Erweiterung der API um
Wetterkorrelation und Sedimentprognosen. Langfristig ist eine Standardisierung des
Datenformats im Rahmen einer IHO-Norm angestrebt.

%=============================================================
\section{Schlussfolgerung}
%=============================================================

Die vorgestellte Lösung zur hochauflösenden Tiefenvermessung adressiert eine bedeutende
Lücke in der maritimen Sicherheitsinfrastruktur. Durch die einmalige Geländeaufnahme
mittels Drohnen oder Satelliten, die Speicherung der prozessierten Daten als ideales
20\,cm-Raster in einer räumlichen Datenbank und die Bereitstellung über eine REST-API
mit einer Latenz unter 500\,ms entsteht eine flexible, skalierbare und kosteneffiziente
Lösung. Die API stellt die Navigationsdaten in Echtzeit auf Grundlage des idealen Rasters
bereit und ermöglicht Schiffen, Booten und autonomen Navigations\-systemen eine
präzise Routenplanung in Flachwasserzonen. Die Ergebnisse des Pilotprojekts belegen
die technische Machbarkeit. Das System eröffnet darüber hinaus neue Perspektiven für
Küstenmanagement, Umweltmonitoring und Such- und Rettungseinsätze.

%=============================================================
% Literaturverzeichnis
%=============================================================

\begin{thebibliography}{99}

% --- Hydrographische Standards ---
\bibitem{iho2020}
International Hydrographic Organization (IHO),
\textit{Standards for Hydrographic Surveys}, 6th Edition,
Special Publication No. 44, Monaco, 2020.

\bibitem{iho2008}
International Hydrographic Organization (IHO),
\textit{IHO Transfer Standard for Digital Hydrographic Data}, Edition 3.1,
Special Publication No. 57, Monaco, 2008.

% --- Maritime Sicherheitsstandards ---
\bibitem{imocircular}
International Maritime Organization (IMO),
\textit{MSC Circular 1023 – Guidelines for Formal Safety Assessment (FSA)},
London, 2002.

\bibitem{imomsccircularnav}
International Maritime Organization (IMO),
\textit{Resolution A.893(21) – Guidelines for Voyage Planning},
London, 1999.

% --- Abtasttheorie ---
\bibitem{shannon1949}
C. E. Shannon,
``Communication in the Presence of Noise,''
\textit{Proceedings of the IRE}, vol.~37, no.~1, pp.~10--21, Jan. 1949.

\bibitem{nyquist1928}
H. Nyquist,
``Certain Topics in Telegraph Transmission Theory,''
\textit{Transactions of the American Institute of Electrical Engineers},
vol.~47, no.~2, pp.~617--644, Apr. 1928.

\bibitem{oppenheim1996}
A. V. Oppenheim und R. W. Schafer,
\textit{Discrete-Time Signal Processing}, 2nd ed.,
Prentice Hall, Upper Saddle River, NJ, 1999.

\bibitem{marks2009}
R. J. Marks II,
\textit{Handbook of Fourier Analysis and Its Applications},
Oxford University Press, New York, 2009.

% --- 2-D Spektraltheorie und Bodenrelief ---
\bibitem{goff1988}
J. A. Goff und T. H. Jordan,
``Stochastic Modeling of Seafloor Morphology: Inversion of Sea Beam Data
for Second-Order Statistics,''
\textit{Journal of Geophysical Research: Solid Earth},
vol.~93, no.~B11, pp.~13589--13608, 1988.

\bibitem{herzfeld1993}
U. C. Herzfeld,
``Least-Squares Collocation, Geophysical Inverse Theory and
Geostatistics: A Bird's Eye View,''
\textit{Geophysical Journal International}, vol.~111, no.~2,
pp.~237--249, 1992.

% --- Rastergeometrien ---
\bibitem{sahr2003}
K. Sahr, D. White und A. J. Kimerling,
``Geodesic Discrete Global Grid Systems,''
\textit{Cartography and Geographic Information Science},
vol.~30, no.~2, pp.~121--134, 2003.

\bibitem{conway2013}
J. H. Conway und N. J. A. Sloane,
\textit{Sphere Packings, Lattices and Groups}, 3rd ed.,
Springer, New York, 1999.

\bibitem{li2008}
Z. Li, Q. Zhu und C. Gold,
\textit{Digital Terrain Modeling: Principles and Methodology},
CRC Press, Boca Raton, FL, 2005.

% --- Bathymetrie und Fernerkundung ---
\bibitem{minkoff2019}
D. Minkoff, J. L. Smith und C. F. Balogh,
``Near-real-time coastal bathymetry from autonomous surface vehicles,''
\textit{Journal of Atmospheric and Oceanic Technology}, vol.~36, no.~8,
pp.~1501--1516, 2019.

\bibitem{guenther2000}
G. C. Guenther, A. G. Cunningham, P. E. LaRocque und D. J. Reid,
``Meeting the accuracy challenge in airborne lidar bathymetry,''
in \textit{Proc. EARSeL Workshop on Lidar Remote Sensing of Land and Sea},
Dresden, Germany, pp.~1--21, 2000.

\bibitem{parente2020}
C. Parente und M. Vallario,
``Interpolation of single beam echo sounder data for 3D coastal bathymetric model,''
\textit{International Journal of Advanced Computer Science and Applications},
vol.~11, no.~11, pp.~575--583, 2020.

\bibitem{caballero2014}
I. Caballero und R. Stumpf,
``Retrieval of nearshore bathymetry from Sentinel-2A and 2B satellites
in South Florida coastal waters,''
\textit{Estuarine, Coastal and Shelf Science}, vol.~214, pp.~189--203, 2019.

\bibitem{collin2018}
A. Collin, B. Long und P. Archambault,
``Merging land-marine realms: Spatial patterns of seamless coastal
topobathymetry,''
\textit{Geomorphology}, vol.~310, pp.~66--74, 2018.

% --- Räumliche Datenbanken und Geodateninfrastruktur ---
\bibitem{argo2021}
M. Argo, T. Schneider und B. Hauschildt,
``Stationary underwater sensor networks for continuous hydrographic monitoring,''
in \textit{Proc. OCEANS 2021: San Diego–Porto}, IEEE, 2021.

\bibitem{stationary2022}
F. Hofmann und D. Kuhlmann,
``Dense fixed-point sensor arrays for coastal depth monitoring,''
\textit{IEEE Journal of Oceanic Engineering}, vol.~47, no.~2,
pp.~411--425, 2022.

\bibitem{obe2015}
R. O. Obe und L. S. Hsu,
\textit{PostGIS in Action}, 2nd ed.,
Manning Publications, Shelter Island, NY, 2015.

\bibitem{ogc2010}
Open Geospatial Consortium (OGC),
\textit{OGC Web Feature Service 2.0 Interface Standard},
OGC Document 09-025r2, 2014.

% --- Datenkompression und Übertragung ---
\bibitem{sayood2017}
K. Sayood,
\textit{Introduction to Data Compression}, 5th ed.,
Morgan Kaufmann, 2017.

\bibitem{wessel2013}
P. Wessel, W. H. F. Smith, R. Scharroo, J. Luis und F. Wobbe,
``Generic Mapping Tools: Improved Version Released,''
\textit{Eos, Transactions American Geophysical Union},
vol.~94, no.~45, pp.~409--410, 2013.

% --- Interpolation und Geostatistik ---
\bibitem{cressie1993}
N. Cressie,
\textit{Statistics for Spatial Data}, revised ed.,
Wiley-Interscience, New York, 1993.


\end{thebibliography}

\end{document}
